\submittedby{Harisankar B}

For $P=(x,y)\in\R^2$ write $\lVert P\rVert=\sqrt{x^2+y^2}$ for its distance from the origin $O$.
Thus $P\sim Q$ if and only if $\lVert P\rVert=\lVert Q\rVert$.

\paragraph{(i)} Since $\sim$ is defined by equality of the real numbers $\lVert\cdot\rVert$, it
inherits the properties of equality:
\begin{itemize}
  \item \emph{Reflexive:} $\lVert P\rVert=\lVert P\rVert$, so $P\sim P$.
  \item \emph{Symmetric:} if $\lVert P\rVert=\lVert Q\rVert$ then $\lVert Q\rVert=\lVert P\rVert$,
        so $P\sim Q\Rightarrow Q\sim P$.
  \item \emph{Transitive:} if $\lVert P\rVert=\lVert Q\rVert$ and $\lVert Q\rVert=\lVert R\rVert$
        then $\lVert P\rVert=\lVert R\rVert$, so $P\sim Q$, $Q\sim R\Rightarrow P\sim R$.
\end{itemize}
Hence $\sim$ is an equivalence relation.

\paragraph{(ii)} Let $r=\lVert P\rVert\ge 0$. Then
\[ [P]=\{Q\in\R^2\mid \lVert Q\rVert=r\}=\{(x,y)\in\R^2\mid x^2+y^2=r^2\}. \]
If $P\neq O$ (so $r>0$), $[P]$ is the circle of radius $r$ centred at the origin passing
through $P$. If $P=O$, then $[P]=\{O\}$, a single point (a ``circle of radius $0$'').

\paragraph{(iii)} The equivalence classes partition $\R^2$ into the origin together with the
family of concentric circles $\{x^2+y^2=r^2\}$, $r>0$, centred at the origin. Every point of
the plane lies on exactly one such circle (the one of radius equal to its distance from $O$),
and distinct radii give disjoint circles. The map $[P]\mapsto\lVert P\rVert$ is a bijection
from the set of equivalence classes $\R^2/{\sim}$ onto $[0,\infty)$, so the classes are
labelled exactly by the possible radii.
